% Generated draft: requires mathematical and source review.
\documentclass[11pt]{article}
\usepackage[margin=25mm]{geometry}
\usepackage{fontspec}
\setmainfont{Latin Modern Roman}
\usepackage{amsmath,amssymb,mathtools,unicode-math}
\setmathfont{Latin Modern Math}
\usepackage{enumitem,needspace,xurl,hyperref}
\setlength{\parindent}{0pt}
\setlength{\parskip}{5pt}
\setlength{\emergencystretch}{3em}
\newcommand{\N}{\mathbb{N}}
\newcommand{\Z}{\mathbb{Z}}
\newcommand{\Q}{\mathbb{Q}}
\newcommand{\R}{\mathbb{R}}
\newcommand{\C}{\mathbb{C}}
\newcommand{\F}{\mathbb{F}}
\newcommand{\A}{\mathbb{A}}
\newcommand{\PP}{\mathbb{P}}
\newcommand{\E}{\mathbb{E}}
\newcommand{\eps}{\varepsilon}
\newcommand{\dd}{\,\mathrm{d}}
\begin{document}
\section*{UnsolvedMath ID: 30001449}
\textbf{Source problem code:} catalogue entry
\subsection{Definitions and mathematical statement}
Let \(p\) be a prime and \(G\) a finite group. Work over an algebraically closed field \(k\) of characteristic \(p\). A \(p\)-block \(B\) of \(G\) is specified by a primitive central idempotent \(e_B\) of \(kG\). Write \(\operatorname{Irr}(B)\) for the ordinary irreducible complex characters belonging to \(B\) under the usual \(p\)-block partition of \(\operatorname{Irr}(G)\).

Let \(D\) be a defect group of \(B\), and write \(|G|_p=p^a\) and \(|D|=p^d\), where \(n_p\) denotes the largest power of \(p\) dividing a positive integer \(n\). For \(\chi\in\operatorname{Irr}(B)\), its height \(h_B(\chi)\) is the nonnegative integer determined by
\[
\chi(1)_p=p^{a-d+h_B(\chi)}.
\]
Define
\[
\operatorname{Irr}_0(B)=\{\chi\in\operatorname{Irr}(B):h_B(\chi)=0\}.
\]
Here \(\chi(1)\) is the full ordinary character degree.

For a \(p\)-subgroup \(Q\leq G\), the Brauer map
\[
\operatorname{Br}_Q:(kG)^Q\longrightarrow kC_G(Q)
\]
retains precisely the coefficients of group elements in \(C_G(Q)\); the superscript denotes invariants under conjugation by \(Q\). A \(B\)-Brauer pair is a pair \((Q,b_Q)\), where \(b_Q\) is a block of \(kC_G(Q)\) with block idempotent \(e_{b_Q}\) satisfying
\[
\operatorname{Br}_Q(e_B)e_{b_Q}\ne 0.
\]
Set
\[
N_G(Q,b_Q)=\{g\in N_G(Q):{}^g b_Q=b_Q\}.
\]
The block \(B\) is called nilpotent if, for every \(B\)-Brauer pair \((Q,b_Q)\), the quotient \(N_G(Q,b_Q)/C_G(Q)\) is a \(p\)-group, allowing the trivial group.

The supplied question asks whether, for every prime \(p\), every finite group \(G\), and every \(p\)-block \(B\) of \(G\),
\[
B\text{ is nilpotent}
\quad\Longleftrightarrow\quad
\forall\chi,\psi\in\operatorname{Irr}_0(B),\qquad \chi(1)=\psi(1).
\]
No restriction to positive defect or to a particular family of finite groups is stated in the supplied formulation.

The record labels this question partially solved and reports results for broad families, including quasi-simple groups, while describing the unrestricted formulation as unresolved. These are supplied dated curation claims, not independently verified conclusions of this draft. The precise hypotheses and scope of the cited partial results are not provided.
\subsection{Short English statement}
For every finite group and every prime \(p\), is a \(p\)-block nilpotent exactly when all its height-zero ordinary irreducible characters have the same full degree? The supplied record describes partial results; its literature assessment requires independent verification.
\subsection{Sources}
\begin{itemize}
\item Local UnsolvedMath registry: supplied complete record for permanent ID \(30001449\), entitled \emph{Character-Degree Criteria for Nilpotent Blocks}; exact source problem code \texttt{catalogue entry}. The record assigns proposed year \(2010\) and supplies no named proposer.
\item Supplied report citation: \emph{Combinatorial Representation Theory} (2010), Oberwolfach Reports, pp.~799--882, DOI \texttt{10.4171/owr/2010/15}. \url{https://doi.org/10.4171/owr/2010/15}. This is also the literature-source entry labeled ``original OWR report.'' No problem-specific page is supplied.
\item Supplied literature source: ``Malle--Navarro, equal height-zero degrees.'' \url{https://arxiv.org/abs/0909.4190}.
\item Supplied literature source: ``Kessar--Malle, quasi-isolated blocks.'' \url{https://annals.math.princeton.edu/2013/178-1/p06}.
\item Literature assessment date: \texttt{2026-08-21}. The supplied dated curation labels the problem \texttt{partially\_solved}, states that the characterization is proved for broad families including quasi-simple groups, and describes the all-finite-groups formulation as a stronger unresolved statement. This assessment and its cited claims have not been independently verified here.
\item The registry assigns statement status \texttt{corrected\_verified} and reports that the mathematical content was checked against the official Oberwolfach report, with notation, context, and question form normalized without changing the source-level claim. This is supplied source-verification metadata, not a verification performed for this draft.
\end{itemize}
\end{document}
